\section{Emergent relativity and effective Lorentz symmetry}
\label{sec:emergent-relativity}
\label{emergent-relativity-and-lorentz-invariance}

A central feature of the brane model is that special-relativistic kinematics
emerges as an effective description of wave propagation in the substrate, rather
than being imposed as a fundamental symmetry. This section outlines how Lorentz
invariance arises in the linear regime and discusses the expected regime of
validity.

\subsection{Wave-cone invariance and the effective light cone}

In the linear wave regime (small amplitudes, weak nonlinearity), the brane
substrate supports propagating modes whose dispersion relation is approximately

\begin{equation}
\omega^2(\mathbf{k}) \approx c^2 |\mathbf{k}|^2,
\label{eq:linear-dispersion}
\end{equation}

where $c = \sqrt{T/\rho_m}$ is the wave speed determined by the brane's tension
$T$ and mass density $\rho_m$. This dispersion relation is isotropic (depends
only on $|\mathbf{k}|$, not on the direction of $\mathbf{k}$) provided the
discrete lattice structure is sufficiently isotropic (e.g., FCC lattice with
appropriate neighbor weights; see Section~\ref{lattice-and-neighborhoods}).

For an observer embedded in the brane who can only access information through
wave propagation, the set of events that can be causally connected to a given
spacetime point $(t_0, \mathbf{x}_0)$ is bounded by the \emph{effective light
cone}:

\begin{equation}
(t - t_0)^2 - \frac{|\mathbf{x} - \mathbf{x}_0|^2}{c^2} = 0.
\label{eq:light-cone}
\end{equation}

Inside this cone (timelike separation), signals can propagate; outside
(spacelike separation), no subluminal signal can connect the events. This
structure is precisely the kinematic foundation of special relativity.

\subsection{Emergent Minkowski metric}

The dispersion relation~\eqref{eq:linear-dispersion} implies that the phase
velocity of a wave packet with central wavevector $\mathbf{k}$ is $\mathbf{v}
= c^2 \mathbf{k} / \omega$, and the group velocity is $\mathbf{v}_g = \partial
\omega / \partial \mathbf{k} = c^2 \mathbf{k} / \omega$. For a wave packet
centered at $\omega_0$, $\mathbf{k}_0$ satisfying the dispersion relation, the
spacetime interval between events separated by $dt$ and $d\mathbf{x}$ is
invariant under transformations that preserve the cone structure. This leads
naturally to the effective Minkowski metric:

\begin{equation}
ds^2 = c^2 dt^2 - d\mathbf{x}^2 = c^2 dt^2 - (dx^1)^2 - (dx^2)^2 - (dx^3)^2.
\label{eq:minkowski-metric}
\end{equation}

Transformations that preserve this metric form the Lorentz group $\mathrm{SO}(1,3)$.

\subsection{Regime of validity and corrections}

The emergent Lorentz invariance is a \emph{regime-based} property, valid in the
limit where:

\begin{enumerate}
\item \textbf{Wavelengths are large compared to the lattice spacing:}
  $|\mathbf{k}| h \ll 1$, where $h$ is the discrete lattice spacing. At shorter
  wavelengths, lattice dispersion corrections become significant and the
  dispersion relation deviates from the linear form~\eqref{eq:linear-dispersion}.

\item \textbf{Amplitudes are small:} Geometric nonlinearity (lateral contraction
  induced by normal displacements) is weak. When amplitudes approach the
  Compton scale $\lambda_C$, the geometric nonlinearity parameter $\eta$
  (Section~\ref{threshold-localization}) becomes of order unity, and the wave
  dynamics can no longer be linearized.

\item \textbf{The medium is isotropic:} The discrete lattice structure and the
  energy functional produce an isotropic effective stiffness tensor. In
  practice, this requires careful choice of lattice geometry and neighbor
  weights (see Section~\ref{calibration-dispersion-isotropy}).
\end{enumerate}

Outside these regimes, corrections to Lorentz invariance are expected. In
particular, at Compton scales (where soliton formation occurs), the effective
description must account for strong geometric nonlinearity and the induced
metric $g_{ij} = \partial_i \mathbf{X} \cdot \partial_j \mathbf{X}$ deviates
from the flat Minkowski form. The framework does not claim exact Lorentz
invariance at all scales; rather, it posits that the observed approximate
Lorentz symmetry of macroscopic physics emerges from the isotropic wave
propagation properties of the substrate in the appropriate regime.

\subsection{Lorentz transformations as effective coordinate changes}

For an internal observer who can only infer spacetime structure through wave
kinematics, a \emph{boost} (change of reference frame with relative velocity
$\mathbf{v}$) corresponds to a re-labeling of events that preserves the causal
structure encoded in the light cone. The standard Lorentz transformation
formulae

\begin{align}
t' &= \gamma \left( t - \frac{\mathbf{v} \cdot \mathbf{x}}{c^2} \right), \\
\mathbf{x}' &= \mathbf{x} + (\gamma - 1) \frac{(\mathbf{v} \cdot \mathbf{x})
\mathbf{v}}{v^2} - \gamma \mathbf{v} t,
\end{align}

with $\gamma = (1 - v^2/c^2)^{-1/2}$, arise as the unique linear transformations
that preserve the interval $ds^2$ and map forward light cones to forward light
cones.

From the brane perspective, these transformations do not reflect a deep symmetry
of the microscopic dynamics (the brane Lagrangian is written with an absolute
time parameter $t$), but rather the kinematic constraints on wave propagation in
an isotropic elastic medium. This is analogous to the emergence of Lorentz-like
symmetry in analog gravity systems (e.g., sound waves in Bose--Einstein
condensates~\cite{Barcelo2011AnalogueGravity}), where the "speed of light" is
replaced by the speed of sound and the symmetry is approximate and regime-dependent.